\section{Explicit Formulas, Positivity Determinants, and Iterated Theta Lifts}
\label{sec:ch02-explicit-formula-theta-lifts}

A classical approach to the Riemann Hypothesis seeks a quadratic form or
determinant whose positivity forces the zeros of Dirichlet $L$-functions onto
the critical line $\operatorname{Re}(s) = \tfrac{1}{2}$.
In function fields over finite fields, Stepanov's method achieves this by
constructing an auxiliary function vanishing to high order $m$ at a large set
of points, thereby converting an analytic zero-counting problem into an
algebraic nonvanishing bound.
Over the integers, however, there is no literal differentiation operator on
prime numbers.
The central analytic-automorphic mechanism of the formalized proof replaces
high-order vanishing with a chain of $m \asymp \sqrt{X}$ iterated theta lifts
between totally definite unitary groups, controlled simultaneously by
Arakelov slope inequalities, determinant positivity bounds, integral dual Weyl
modules, and Moriwaki arithmetic bigness.

In this chapter we expose this analytic engine directly from the verified Lean
development, progressing from the critical-line reduction and Selmer-window
formulation through Euclidean and Hermitian lattice slopes, iterated definite
theta lifts, graded Weyl divisibility, and the arithmetic height bounds that
close the global inequality.

% -----------------------------------------------------------------------------
\subsection{From Half-Plane Nonvanishing to the Critical Line}
\label{subsec:ch02-critical-line}

The global analytic target of the proof is the universal right-half-plane
nonvanishing statement \leanlink{UpperTheory.EveryCharacter}: for every modulus
$N \ge 1$, every Dirichlet character $\chi \pmod{N}$, and every
$s \in \mathbb{C}$ with $\operatorname{Re}(s) > \tfrac{1}{2}$ (and
either $\chi \neq 1$ or $s \neq 1$), one has $L(s, \chi) \neq 0$.
Once this half-plane nonvanishing is established, the functional equation of
completed Dirichlet $L$-functions reflects the zero-free region across
$\operatorname{Re}(s) = \tfrac{1}{2}$ and confines every nontrivial zero to the
critical line.

The reflection argument relies on two elementary but essential analytic facts:
the conductor of a Dirichlet character is invariant under inversion, and the
archimedean gamma factor has no zeros in the right half-plane
$\operatorname{Re}(s) > 0$.

\begin{lemma}[Conductor inversion and gamma-factor nonvanishing]
\label{lem:ch02-gamma-nonvanishing}
\leanlink{Formalism.VW.gammaFactor_ne_zero}
Let $N \in \mathbb{N}$, let $\chi$ be a Dirichlet character modulo $N$, and let
$s \in \mathbb{C}$ satisfy $0 < \operatorname{Re}(s)$.
Then the archimedean gamma factor satisfies $\gamma(s, \chi) \neq 0$.
Moreover, the conductor of $\chi^{-1}$ coincides with the conductor of $\chi$
\leanlink{PI.EF.conductor_inv}.
\end{lemma}

\begin{proof}
Every Dirichlet character $\chi$ is either even or odd.
In the even case $\gamma(s, \chi) = \Gamma_{\mathbb{R}}(s) = \pi^{-s/2}\Gamma(s/2)$,
and in the odd case $\gamma(s, \chi) = \Gamma_{\mathbb{R}}(s + 1)$.
Since $\Gamma_{\mathbb{R}}(z) \neq 0$ whenever $\operatorname{Re}(z) > 0$,
the hypothesis $0 < \operatorname{Re}(s)$ guarantees $\gamma(s, \chi) \neq 0$ in
both cases.
\end{proof}

For the Riemann zeta function itself, right-half-plane nonvanishing immediately
yields the classical Riemann Hypothesis via Riemann's functional equation.

\begin{proposition}[Critical-line reflection for $\zeta(s)$]
\label{prop:ch02-rh-of-right}
\leanlink{Formalism.VW.rh_of_right}
Suppose that for every $s \in \mathbb{C}$ satisfying $\tfrac{1}{2} < \operatorname{Re}(s)$
and $s \neq 1$, we have $\zeta(s) \neq 0$.
Then the Riemann Hypothesis holds: every zero $s \in \mathbb{C}$ of $\zeta(s)$
other than the trivial zeros $s = -2k$ ($k \in \mathbb{Z}_{\ge 1}$) and $s = 1$
satisfies $\operatorname{Re}(s) = \tfrac{1}{2}$.
\end{proposition}

\begin{proof}
Let $s \in \mathbb{C}$ be a nontrivial zero of $\zeta(s)$ with $s \neq 1$.
By hypothesis, $\operatorname{Re}(s)$ cannot strictly exceed $\tfrac{1}{2}$.
Suppose for contradiction that $\operatorname{Re}(s) < \tfrac{1}{2}$.
Since $\zeta(0) = -\tfrac{1}{2} \neq 0$, we have $s \neq 0$, so
$w := 1 - s$ satisfies $\operatorname{Re}(w) > \tfrac{1}{2}$ and $w \neq 1$,
and in particular $w \notin \mathbb{Z}_{\le 0}$.
By Riemann's functional equation,
\[
  \zeta(s)
  = 2 (2\pi)^{-w}\,\Gamma(w)\,\cos\!\left(\frac{\pi w}{2}\right)\,\zeta(w).
\]
On the right-hand side, $2(2\pi)^{-w} \neq 0$, $\Gamma(w) \neq 0$ because
$\operatorname{Re}(w) > 0$, and $\zeta(w) \neq 0$ by the right-half-plane
assumption.
Hence $\cos(\pi w / 2) = 0$, which forces $w = 2k + 1$, i.e., $s = -2k$, for
some integer $k$.
The inequalities $\operatorname{Re}(s) < \tfrac{1}{2}$ and $s \neq 0$ rule out
$k \le 0$, leaving $k \ge 1$ and contradicting the nontriviality of $s$.
Therefore $\operatorname{Re}(s) = \tfrac{1}{2}$.
\end{proof}

Applying the same reflection mechanism to the completed $L$-function
$\Lambda(s, \chi) = \gamma(s, \chi)\,L(s, \chi)$ of an arbitrary primitive
Dirichlet character gives both the Generalized Riemann Hypothesis on the
critical strip and the classical Riemann Hypothesis as direct consequences of
\leanlink{UpperTheory.EveryCharacter}.

\begin{theorem}[Critical-line localization for primitive Dirichlet $L$-functions]
\label{thm:ch02-crit-line}
\leanlink{Formalism.VW.critLine}
Assume \leanlink{UpperTheory.EveryCharacter}.
Let $N \ge 1$, let $\chi$ be a primitive Dirichlet character modulo $N$, and
let $s \in \mathbb{C}$ satisfy $0 < \operatorname{Re}(s) < 1$ and
$L(s, \chi) = 0$.
Then $\operatorname{Re}(s) = \tfrac{1}{2}$.
\end{theorem}

\begin{proof}
Because $0 < \operatorname{Re}(s) < 1$, we have $s \neq 0$ and $s \neq 1$.
The assumption \leanlink{UpperTheory.EveryCharacter} immediately rules out
$\operatorname{Re}(s) > \tfrac{1}{2}$.
If $\operatorname{Re}(s) < \tfrac{1}{2}$, then $\chi^{-1}$ is again primitive
by conductor invariance \leanlink{PI.EF.conductor_inv}.
Writing $L(s, \chi) = \Lambda(s, \chi) / \gamma(s, \chi)$ and applying
Lemma~\ref{lem:ch02-gamma-nonvanishing} (\leanlink{Formalism.VW.gammaFactor_ne_zero})
gives $\Lambda(s, \chi) = 0$.
By the functional equation relating $\Lambda(1 - s, \chi^{-1})$ and
$\Lambda(s, \chi)$, it follows that $\Lambda(1 - s, \chi^{-1}) = 0$, and hence
$L(1 - s, \chi^{-1}) = 0$ at $1 - s$ with $\operatorname{Re}(1 - s) > \tfrac{1}{2}$
and $1 - s \neq 1$, contradicting \leanlink{UpperTheory.EveryCharacter}.
\end{proof}

\begin{corollary}[Riemann Hypothesis from universal character nonvanishing]
\label{cor:ch02-rh-of-every}
\leanlink{Formalism.VW.rh_of_every}
If \leanlink{UpperTheory.EveryCharacter} holds, then the Riemann Hypothesis
holds.
\end{corollary}

\begin{proof}
Specializing \leanlink{UpperTheory.EveryCharacter} to the trivial character of
modulus $N = 1$, whose $L$-function equals $\zeta(s)$, supplies the hypothesis
of Proposition~\ref{prop:ch02-rh-of-right} (\leanlink{Formalism.VW.rh_of_right}).
\end{proof}

To prove \leanlink{UpperTheory.EveryCharacter}, the formalization factors the
analytic explicit-formula contradiction through a quantitative Vandiver-type
cohomological window condition \leanlink{UpperTheory.Top.VWin} (defined in
\leanmod{Development.AutomorphicArithmetic.CharacterBoxes.ExistsPageN}).
For each primitive character $\psi$, small exponent $0 < \varepsilon < 1/52$,
and scale $x \ge 1$, one selects a dyadic level $2^e \in [x^{2 + 8\varepsilon},
2 x^{2 + 8\varepsilon}]$ (\leanlink{UpperTheory.Top.PageN}) and a pair of
auxiliary primes $(q, d)$ satisfying the congruence and logarithmic-size
constraints \leanlink{UpperTheory.Top.WinPair}:
\[
  q \equiv 3 \pmod{8}, \qquad d \equiv 7 \pmod{16}, \qquad q \mid d + 1,
\]
with $\log q \asymp Y^{\varepsilon/5}$,
$\log d \asymp Y^{\varepsilon/5}(\log Y)^2$, and $p \mid d + 1$ for every odd
prime $p \le Y^{\varepsilon/5}$.
To each tuple $(\psi, e, q, d)$ is attached an abelian number field $E$
(\leanlink{UpperTheory.Top.IsChiField}) cut out inside a cyclotomic field
$\mathbb{Q}(\zeta_N)$ by the subgroup where $a \equiv \pm 1 \pmod{2^{e+2}}$,
$\psi(a) = 1$, and $\bigl(\frac{a}{q}\bigr) = \bigl(\frac{a}{d}\bigr) = 1$.
The window predicate \leanlink{UpperTheory.Top.VAt} requires that for every
finite set $S$ of primes, the Tate--Shafarevich / étale cohomology length
weighted by $\log \ell$ obeys the subexponential bound
\leanlink{UpperTheory.Top.VVle}:
\[
  \sum_{\ell \in S} (\log \ell) \sum_{1 \le k < \lfloor x^{401} \rfloor}
    \operatorname{len}\!\left(\operatorname{Sel}_\ell(E, k)\right)
  \;\le\;
  \exp\!\left(\tfrac{1}{2} x^{\varepsilon}\right).
\]
Assembling the automorphic window theorem \leanlink{UpperTheory.One.vwin_of_row}
with the analytic-arithmetic implication
\leanlink{UpperTheory.LiteratureTheorems.every_character_of_literature_of_vWin}
closes the top-level deduction of \leanlink{UpperTheory.EveryCharacter} from
the Galois-automorphic row hypothesis \leanlink{UpperTheory.GalStar.Row0b}
via \leanlink{UpperTheory.FinalTheorems.everyCharacter_of_row}.

% -----------------------------------------------------------------------------
\subsection{Stepanov-Style Multiplicity, Arakelov Slopes, and Positivity Determinants}
\label{subsec:ch02-slopes-determinants}

Why does a bound on Selmer cohomology lengths over cyclotomic-quadratic fields
prevent an off-line zero of $L(s, \chi)$?
Conceptually, the Guinand--Weil explicit formula converts a hypothetical zero
with $\operatorname{Re}(\rho) > \tfrac{1}{2}$ into a persistent macroscopic
imbalance in weighted von Mangoldt sums $\sum_{n \le X} \Lambda(n)\chi(n) w(n)$
across dyadic windows.
In the proof's quantitative core (\leanmod{Formalism.SizeWindowBrackets}),
these size windows are parameterized by the structure
\leanlink{Formalism.LFunctionChainParameters.PrimeWindowSizeBounds.Par}
together with the explicit threshold scaling
\leanlink{Formalism.LFunctionChainParameters.PrimeWindowSizeBounds.Facts.hsg},
which evaluates the auxiliary window parameter as
$\operatorname{sg}(P, K) = 100 K z^4 P.x^{-1}$.
To contradict this imbalance in the spirit of Stepanov, one constructs a
lattice of integral automorphic sections whose evaluation map has small
operator norm at the archimedean places while possessing high divisibility at
finite places.
The comparison between the finite-place divisibility gain and the archimedean
height loss is governed by Bost's Arakelov slope inequality for Euclidean and
Hermitian lattices.

In \leanmod{Formalism.LatticeSlopes} and \leanmod{Formalism.HarderNarasimhan},
the formalization builds this slope calculus from first principles.
For a discrete $\mathbb{Z}$-submodule $L$ of a finite-dimensional real inner
product space $E$, its covolume (\leanlink{Formalism.F.covol}), rank
(\leanlink{Formalism.F.rk}), Arakelov degree (\leanlink{Formalism.F.deg}), and
slope (\leanlink{Formalism.F.slope}) are defined by
\[
  \operatorname{covol}(L) = \sqrt{\det\bigl(\operatorname{Gram}(b_1, \dots, b_r)\bigr)},
  \qquad
  \widehat{\deg}(L) = -\log\operatorname{covol}(L),
  \qquad
  \mu(L) = \frac{\widehat{\deg}(L)}{\operatorname{rk}(L)},
\]
where $(b_1, \dots, b_r)$ is any $\mathbb{Z}$-basis of $L$.
For a sublattice $N \le L$, its saturation
$\operatorname{satur}_L(N) = L \cap \operatorname{span}_{\mathbb{R}}(N)$
maximizes the Arakelov degree among all sublattices of $L$ spanning the same
real subspace, because passing to a finite-index superlattice divides the
covolume by the index and strictly increases $-\log\operatorname{covol}$.

\begin{proposition}[Lattice saturation and Harder--Narasimhan rank identity]
\label{prop:ch02-lattice-saturation}
\leanlink{Formalism.F.saturation_properties}
Let $E$ be a finite-dimensional real inner product space and $L \subset E$ a
discrete $\mathbb{Z}$-lattice.
For any sublattice $N \le L$, the saturation $\operatorname{satur}_L(N)$ is
saturated in $L$, contains $N$, and satisfies
$\operatorname{rk}(\operatorname{satur}_L(N)) = \operatorname{rk}(N)$.
In particular, for every discrete lattice $N \subset E$, the algebraic rank
equals the dimension of its real span \leanlink{Formalism.F.rank_eq_finrank_span}:
\[
  \operatorname{rk}(N) = \dim_{\mathbb{R}} \operatorname{span}_{\mathbb{R}}(N).
\]
\end{proposition}

\begin{proof}
By \leanlink{Formalism.F.rank_eq_finrank_span}, any discrete
$\mathbb{Z}$-submodule $M \subset E$ becomes a full $\mathbb{Z}$-lattice inside
its own real span $V = \operatorname{span}_{\mathbb{R}}(M)$, so its
$\mathbb{Z}$-rank agrees with $\dim_{\mathbb{R}}(V)$.
Since $N \le \operatorname{satur}_L(N) \le \operatorname{span}_{\mathbb{R}}(N)$,
both $N$ and $\operatorname{satur}_L(N)$ have the same real span, whence
$\operatorname{rk}(\operatorname{satur}_L(N)) = \operatorname{rk}(N)$, and
saturation in $L$ follows immediately from intersecting $L$ with that span.
\end{proof}

Over a number field $K$, the same saturation mechanism extends to
$\mathcal{O}_K$-lattices equipped with Hermitian metrics at all infinite places
$v \mid \infty$ (\leanlink{Formalism.HN.ht}, \leanlink{Formalism.HN.muDef}),
via the real saturation theorem \leanlink{Formalism.HN.real_saturation_property},
which enables the Harder--Narasimhan filtration and Bost's slope inequality
comparing maximal slopes under linear maps of bounded operator norm
(\leanlink{Formalism.HN.opNormAt}).

To evaluate the Arakelov degrees of lattices arising from Fock spaces and
theta series, the formalization proves the exact duality and index-determinant
identities for Hermitian lattices in \leanmod{Formalism.FockGram}.
Given a nondegenerate Hermitian form $B$ on a $K$-vector space $V$ and an
$\mathcal{O}_K$-lattice $L \subset V$, its Hermitian dual lattice
$L^{\#} = \{x \in V : B(x, L) \subseteq \mathcal{O}_K\}$
(\leanlink{Formalism.E.hdual}) satisfies the covolume-Gram and index-norm
relations below.

\begin{proposition}[Hermitian dual lattices and Fock-Gram covolumes]
\label{prop:ch02-fock-gram}
\leanlink{Formalism.E.hermDual}
For finite-rank Hermitian lattices equipped with Euclidean/Hermitian inner
products:
\begin{enumerate}
  \item The covolume of a lattice equals the square root of the determinant of
        its Gram matrix \leanlink{Formalism.E.covolumeEqSqrtDetGram}.
  \item For any full-rank sublattice $M \le L$, the index $[L : M]$ equals the
        absolute norm of the determinant of the change-of-basis matrix
        \leanlink{Formalism.E.subIndexNorm}.
  \item The Hermitian dual lattice $L^{\#}$ satisfies the standard dual-volume
        and duality identities (\leanlink{Formalism.E.hermDual}, proved from
        \leanlink{Formalism.E.subIndexNorm} via \leanlink{Formalism.E.hermDual_of}).
\end{enumerate}
\end{proposition}

Alongside lattice covolumes, the positivity arguments repeatedly compare
determinants of positive semidefinite Gram matrices when passing to subfamilies
or adding positive spectral contributions.
This monotonicity is established directly in
\leanmod{Development.AutomorphicSpectrum.Proofs.BlockDecomposition.DeterminantMonotonicity}.

\begin{proposition}[Hermitian determinant monotonicity]
\label{prop:ch02-det-monotone}
\leanlink{GlobalTheory.CompanionTheory.DetMonotone.detMono}
If $X$ and $Y$ are finite-dimensional complex Hermitian matrices such that $X$
and $Y - X$ are positive semidefinite, then
\[
  \operatorname{Re}(\det X) \;\le\; \operatorname{Re}(\det Y).
\]
\end{proposition}

\begin{proof}
If $\det X = 0$, the inequality reduces to $\operatorname{Re}(\det Y) \ge 0$,
which holds because $Y = X + (Y - X)$ is positive semidefinite.
Otherwise $X$ is strictly positive definite, and unitary diagonalization
(\leanlink{GlobalTheory.CompanionTheory.DetMonotone.exists_unitary_diagonal_eigenvalues_of_isHermitian})
yields $X = B^{*} B$ with $B$ invertible.
Writing $Y = B^{*}(I + Q)B$ where $Q = (B^{*})^{-1}(Y - X)B^{-1}$ is positive
semidefinite reduces the inequality to $\det(I + Q) \ge 1$, which follows from
expanding $\det(I + Q) = \prod_i (1 + \lambda_i(Q))$ with all eigenvalues
$\lambda_i(Q) \ge 0$.
\end{proof}

% -----------------------------------------------------------------------------
\subsection{Iterated Theta Lifts and Definite Unitary Period Integrals}
\label{subsec:ch02-iterated-theta-lifts}

The mechanism that generates effective multiplicity $m$ over $\mathbb{Z}$ is
an $m$-fold chain of theta lifts
\[
  \mathrm{U}(X_0) \longrightarrow \mathrm{U}(X_1) \longrightarrow \cdots
  \longrightarrow \mathrm{U}(X_m)
\]
between totally definite unitary groups.
At each step, lifting a form along the reductive dual pair shifts the
automorphic representation by a Hecke character packet
$\pi \mapsto \pi \oplus \mu_j[t]$, gaining one order of algebraic vanishing
in the Stepanov determinant while multiplying the integral denominators by a
controlled local mass factor.
Because every group $\mathrm{U}(X_j)$ is totally definite, its automorphic
quotient at fixed level decomposes into a finite set $\iota_j$ of isometry
classes of Hermitian lattices (\leanlink{Formalism.ThetaLiftIndependence.IsClassReps}),
and each theta lift acts as an explicit weighted finite-sum kernel operator.

In \leanmod{Formalism.ThetaLiftIndependence.Statement}, the formalization
captures this discrete kernel algebra cleanly over any commutative ring $R$
and any sequence of finite class sets $(\iota_j)_{j \ge 0}$.
Given weights $w_j : \iota_j \to R$ (encoding inverse automorphism-group sizes
$|\operatorname{Aut}(L_i)|^{-1}$, with lattice automorphism groups formalized
by \leanlink{Formalism.ThetaLiftIndependence.autL}) and module-valued kernel
operators $K_j(i, i') : M_j \to M_{j+1}$, the one-step theta lift operator
\leanlink{Formalism.ThetaLiftIndependence.kerOpM} is
\[
  (\mathcal{K}_{w_j, K_j} f)(y) = \sum_{x \in \iota_j} w_j(x)\, K_j(x, y)\bigl(f(x)\bigr),
\]
and the $m$-fold iterated lift \leanlink{Formalism.ThetaLiftIndependence.chainOpM}
is defined inductively by
\[
  \mathcal{C}_0(f) = f,
  \qquad
  \mathcal{C}_{m+1}(f) = \mathcal{K}_{w_m, K_m}\!\bigl(\mathcal{C}_m(f)\bigr).
\]
This formulation makes three structural properties transparent:
\begin{itemize}
  \item \emph{Adjointness, composition, and multiplicative denominator control}:
        with respect to the mass-weighted inner product
        $\langle f, g \rangle_w = \sum_i w_i \overline{f(i)} g(i)$, the adjoint
        of a kernel operator is given by the conjugate-transpose kernel, and
        if scalars $D_j \in R$ clear the mass denominators $D_j w_j(x)$ while
        each kernel $K_j(x, y)$ maps an integral lattice $\Lambda_j \subset M_j$
        into $\Lambda_{j+1} \subset M_{j+1}$, then the $m$-fold iterated lift
        satisfies
        \[
          \left(\prod_{j=0}^{m-1} D_j\right) \cdot \mathcal{C}_m(f)(y) \in \Lambda_m
        \]
        for every integral input $f$ with $f(x) \in \Lambda_0$.
  \item \emph{Definite class-set independence}: for a definite Hermitian space
        $(V, h)$, the action of isometries on $\mathcal{O}_K$-lattices is
        functorial (\leanlink{Formalism.ThetaLiftIndependence.LatticeLiftIndependence.gmap_mul},
        \leanlink{Formalism.ThetaLiftIndependence.LatticeLiftIndependence.gmap_inv_gmap})
        and induces an isometry equivalence relation
        (\leanlink{Formalism.ThetaLiftIndependence.DefiniteClassRepresentatives.isoSetoid})
        whose chosen class representatives
        (\leanlink{Formalism.ThetaLiftIndependence.DefiniteClassRepresentatives.pick})
        preserve membership in finite genus sets
        (\leanlink{Formalism.ThetaLiftIndependence.DefiniteClassRepresentatives.pick_spec}).
  \item \emph{Finite-field Plancherel orthogonality}: local character-sum
        cancellation and orthogonality over finite residue fields are supplied
        by \leanlink{Formalism.ThetaLiftIndependence.FiniteFieldPlancherel.charSum}
        and \leanlink{Formalism.ThetaLiftIndependence.FiniteFieldPlancherel.orth}.
\end{itemize}

To evaluate the norms and periods of these iterated theta lifts, one invokes
the seesaw identity and the Rallis inner product formula, which express the
period of an iterated lift in terms of Bessel periods and local formal degrees.
Two key pillars of this local-global period machinery appear in
\leanmod{Formalism.PeriodIntegrals} and \leanmod{Formalism.FormalDegreeBessel}.

First, on level-zero Hermitian matrix spaces equipped with a Haar measure
$\mu$ satisfying the smooth-test-space axioms \leanlink{PI.QnHermHaar.Hyp},
the differential / Lie-algebra operators $E_{a,b}$ acting on the test space
\leanlink{PI.QnHermHaar.testSub} admit explicit formal adjoints $E'_{a,b}$
under integration against $\mu$.

\begin{proposition}[Formal adjointness on level-zero Hermitian Haar spaces]
\label{prop:ch02-herm-haar-adj}
\leanlink{PI.QnHermHaar.adjOp_E}
Let $k, r \in \mathbb{N}$, $i \in \{0, \dots, r\}$, and let $\mu$ be a measure
on the Hermitian matrix space satisfying the level-zero hypothesis
\leanlink{PI.QnHermHaar.Hyp}.
Then for every pair of matrix indices $a, b$, the operator $E_{a,b}$ on the
test submodule $\mathcal{F} = \operatorname{testSub}_{k,r,i}$
(\leanlink{PI.QnHermHaar.testSub}) has formal $\mu$-adjoint $E'_{a,b}$
relative to the Hermitian pairing.
\end{proposition}

\begin{proof}
This follows by integrating by parts on the level-zero Hermitian chart and
identifying the transposed action via \leanlink{PI.QnHermHaar.EOp'_eq}.
\end{proof}

Second, the archimedean formal-degree ratios and Bessel positivity estimates
rely on algebraic identities for quantum-determinant projectors $H_{p,k}$
acting on Hermitian spaces of signature $(p, q)$.

\begin{theorem}[Hermitian positivity identities for Bessel formal-degree projectors]
\label{thm:ch02-bessel-pos-facts}
\leanlink{Formalism.QuantumDeterminants.hermitian_pos_facts_hold}
The projector family $H_{p,k}$ on Hermitian representation spaces satisfies the
complete package of positivity, idempotence, and symmetry relations
(exported as \leanlink{Formalism.QuantumDeterminants.hpkFacts}): each $H_{p,k}$
is Hermitian, satisfies the quadratic relation $H_{p,k}^2 = H_{p,k}$ up to
explicit scalar normalization, and intertwines the signature form $J$ via the
projected operator identity $P J P$.
\end{theorem}

\begin{proof}
In the formalization, the three structural properties are established
individually as the Hermiticity relation
\leanlink{Formalism.QuantumDeterminants.GBasic.Hpk_herm}, the quadratic
identity \leanlink{Formalism.QuantumDeterminants.GBasic.Hpk_sq}, and the
projected signature identity \leanlink{Formalism.QuantumDeterminants.GBasic.PJP},
which combine directly to yield
\leanlink{Formalism.QuantumDeterminants.hermitian_pos_facts_hold}.
\end{proof}

% -----------------------------------------------------------------------------
\subsection{Graded Weyl Divisibility and Integral Dual Weyl Coordinates}
\label{subsec:ch02-weyl-divisibility}

A critical quantitative hurdle in the iterated theta-lift chain is to prove
sufficiently strong finite-place divisibility for the $m$-fold lift so that
its Arakelov slope beats the archimedean norm growth.
Because the archimedean components of the theta lifts live in algebraic
representations of $\mathrm{GL}_r$ realized inside polynomial Fock spaces
$R[X_{a,b}]_{1 \le a \le r,\, 1 \le b \le n}$, establishing exact divisibility
at primes $p$ requires working over $\mathbb{Z}$ (and arbitrary integral
domains $R$), not merely over $\mathbb{C}$.
This integral representation theory is developed in
\leanmod{Formalism.GradedWeylDivisibility} and
\leanmod{Formalism.ThetaLogCoordinates}.

Over a commutative ring $R$, let $P_{r,n}(R) = R[X_{a,b}]$ be the polynomial
ring in an $r \times n$ matrix of indeterminates.
Given a column-degree tuple $\nu = (\nu_1, \dots, \nu_n) \in \mathbb{N}^n$,
the \emph{integral dual Weyl module}
\leanlink{Formalism.GradedWeylDivisibility.dualWeyl}, denoted
$\nabla_R(r, n; \nu)$, is defined as the intersection of the
$\nu$-weighted column-homogeneous submodule with the common kernel of the
strictly upper-triangular column-addition derivations
\leanlink{Formalism.GradedWeylDivisibility.colAddGen}:
\[
  \nabla_R(r, n; \nu)
  \;=\;
  P_{r,n}(R)_{\nu} \;\cap\;
  \bigcap_{1 \le l < j \le n}
  \ker\bigl(\operatorname{colAddGen}_{l,j} - \iota\bigr).
\]
Equivalently, $\nabla_R(r, n; \nu)$ consists of column-homogeneous polynomials
invariant under the unipotent subgroup $U_n^+$ acting on columns.
Its canonical generators are \emph{bideterminants}
\leanlink{Formalism.GradedWeylDivisibility.bideterminant}: products
$\prod_{k=1}^K \det\bigl(X_{I_k(p), q}\bigr)_{1 \le p, q \le s_k}$ of initial
minors (\leanlink{Formalism.GradedWeylDivisibility.initMinor}) of sizes $s_k$
determined by the conjugate partition of $\nu$.
Row operations $\mathrm{GL}_r(R)$ act on the left via
\leanlink{Formalism.GradedWeylDivisibility.rowAct}, together with divided-power
raising and lowering operators \leanlink{Formalism.GradedWeylDivisibility.rowDiv}.

\begin{theorem}[Highest-weight structure, semistandard peeling, and rational Weyl identification]
\label{thm:ch02-weyl-structure}
\leanlink{Formalism.GradedWeylDivisibility.weylRational}
Let $r, n \in \mathbb{N}$ with $n \le r$ and let $\nu : \{1, \dots, n\} \to \mathbb{N}$
be an antitone weight (a partition).
\begin{enumerate}
  \item \emph{Highest-weight line} (\leanlink{Formalism.GradedWeylDivisibility.highest_weight_line}):
        the highest-weight space inside $\nabla_{\mathbb{Q}}(r, n; \nu)$ is
        one-dimensional, spanned by the canonical product of leading principal
        minors.
  \item \emph{Semistandard tableau peeling} (\leanlink{Formalism.GradedWeylDivisibility.semistandardPeel}):
        iterated leading-monomial extraction peels any nonzero polynomial in
        $\nabla_R(r, n; \nu)$ into semistandard Young tableau bideterminants,
        controlling leading exponents in column-lexicographic order.
  \item \emph{Rational identification} (\leanlink{Formalism.GradedWeylDivisibility.weylRational}):
        over $\mathbb{Q}$ (and hence over $\mathbb{C}$), the Weyl module
        generated by divided-power row operators from the highest-weight line
        coincides with the dual Weyl module $\nabla_{\mathbb{Q}}(r, n; \nu)$.
        Over $\mathbb{C}$, the inclusion
        $\nabla_{\mathbb{C}}(r, n; \nu) \subseteq \Delta_{\mathbb{C}}(r, n; \nu)$
        holds for every antitone $\nu$
        \leanlink{Formalism.BranchingDeterminantCap.bs_dual_le_weyl}.
\end{enumerate}
\end{theorem}

\begin{proof}
In the Lean formalization, the rational identification
\leanlink{Formalism.GradedWeylDivisibility.weylRational} is assembled directly
from \leanlink{Formalism.GradedWeylDivisibility.weylEqNablaQ} and
\leanlink{Formalism.GradedWeylDivisibility.weyl_module_rational}.
Conceptually, the Fischer inner product on the rational Fock space is positive
definite and makes the divided-power row operators adjoint to one another;
combined with uniqueness of the highest-weight line
(\leanlink{Formalism.GradedWeylDivisibility.highest_weight_line}), any
orthogonal complement to the Weyl module inside $\nabla_{\mathbb{Q}}(r, n; \nu)$
would contain a nonzero highest-weight vector and hence vanish.
\end{proof}

To pass from rational representations to integral coordinates on the algebraic
group $\mathrm{GL}_n$ (where negative determinant powers also occur for general
dominant weights $\mu \in \mathbb{Z}^n$),
\leanmod{Formalism.ThetaLogCoordinates.DualWeylCoordinates} proves the integral
straightening law and the standard bideterminant basis theorem over arbitrary
integral domains.

\begin{theorem}[Integral straightening and standard basis theorem for dual Weyl modules]
\label{thm:ch02-standard-basis}
\leanlink{Formalism.InvariantFactorPaths.ColumnFamilies.dualWeyl_eq_span_standard_tableauBideterminant}
Let $B$ be a commutative integral domain, let $n \ge 1$, and let
$\mu \in \mathbb{Z}^n$ be a dominant weight ($\mu_1 \ge \mu_2 \ge \dots \ge \mu_n$).
Then the integral dual Weyl module $\nabla_B(n, n; \nu(\mu))$ is freely
generated over $B$ by the bideterminants indexed by standard tableaux
$I \in \operatorname{StdT}(\mu)$:
\[
  \nabla_B\bigl(n, n; \nu(\mu)\bigr)
  \;=\;
  \operatorname{span}_B\!\bigl\{\, P_B(I) \;\big|\; I \in \operatorname{StdT}(\mu) \,\bigr\}.
\]
Twisting by $(\det U)^{\mu_n}$ inside the coordinate ring of $\mathrm{GL}_n$
over $\mathbb{Z}$ yields the integral standard basis theorem
\leanlink{Formalism.InvariantFactorPaths.ColumnFamilies.standard_basis_theorem}
and the straightening law
\leanlink{Formalism.InvariantFactorPaths.ColumnFamilies.claimA_straightening}.
\end{theorem}

\begin{proof}
Each standard tableau bideterminant lies in the dual Weyl module by
\leanlink{Formalism.InvariantFactorPaths.ColumnFamilies.tableauBideterminant_mem_dualWeyl}.
Conversely, given any element of $\nabla_B(n, n; \nu(\mu))$, its leading
exponent is interlaced by \leanlink{Formalism.GradedWeylDivisibility.leadInterlaced}
and is hit by a standard tableau bideterminant via
\leanlink{Formalism.InvariantFactorPaths.ColumnFamilies.qa_exp_surj}.
Subtracting that scalar multiple strictly lowers the leading term in the
monomial order, so well-founded induction expresses every element as a
$B$-linear combination of standard bideterminants.
\end{proof}

Using the straightening law and standard basis theorem, the formalization
establishes the closure of column-family modules under trivial representations,
dual Weyl modules, and determinant twists
(\leanlink{Formalism.InvariantFactorPaths.ColumnFamilies.closure_first_part_of}),
as well as the generic-point filtration criterion.

\begin{theorem}[Generic-point criterion for good filtrations]
\label{thm:ch02-filtration-criterion}
\leanlink{Formalism.InvariantFactorPaths.ColumnFamilies.filtration_criterion_of}
Assuming the integral straightening law ($\mathrm{LemmaA}$,
\leanlink{Formalism.InvariantFactorPaths.ColumnFamilies.claimA_straightening}),
let $M$ be a rational $\mathrm{GL}_n$-module equipped with a column family
$\mathcal{CF}$, let $\pi \subset \mathbb{Z}^n$ be a lower set of dominant
weights, and let $F$ be any field.
Then a vector $x \in M \otimes_{\mathbb{Z}} F$ lies in the $F$-span of the
standard-tableau basis vectors of weights in $\pi$ if and only if its
generic-point evaluation projections vanish for all dominant weights
$\mu \notin \pi$.
\end{theorem}

\begin{proof}
This is a direct reformulation of
\leanlink{Formalism.InvariantFactorPaths.ColumnFamilies.mem_span_standard_iff_proj_zero},
which combines straightening with the linear independence of standard tableau
evaluations over arbitrary fields $F$.
\end{proof}

These integral dual Weyl coordinates are precisely what permit integral
structures on automorphic representations to be glued across varying Hecke
characters and pushed through the $m$-fold theta-lift chain without losing
control of prime-power denominators.

% -----------------------------------------------------------------------------
\subsection{Arithmetic Heights, Moriwaki Bigness, and Automorphic Positivity}
\label{subsec:ch02-heights-positivity}

The final analytic-geometric link in Chapter~2 connects the algebraic
divisibility of iterated theta lifts to global arithmetic height lower bounds
and automorphic positivity inputs.
Across \leanmod{GroundTheory.HeightsAndTheta},
\leanmod{PositivityTheory}, and \leanmod{Development.AutomorphicSpectrum}, the
Lean proof formalizes the full Arakelov-geometric infrastructure needed to
extract nonzero small integral sections (arithmetic Hilbert--Samuel and
Moriwaki bigness) and to supply automorphic positivity and base-change data.

On the arithmetic-height side (\leanmod{GroundTheory.HeightsAndTheta}),
measuring the arithmetic volume of line bundles on arithmetic varieties
requires relating the transcendence degree of multihomogeneous coordinate cones
to Krull dimension and computing local intersection lengths on integral models:
\begin{itemize}
  \item \emph{Multihomogeneous cone transcendence degree}
        (\leanlink{GroundTheory.Heights₂.Heights.ConeTrdeg.proof}): controls
        the transcendence degree of the multihomogeneous coordinate ring of a
        projective arithmetic subvariety in terms of its geometric dimension.
  \item \emph{Function-field dimension lower bound}
        (\leanlink{GroundTheory.Heights₂.Heights.LowerBoundFree.FnFieldTrdeg.full}):
        establishes the transcendence-degree and Krull-dimension bridge used in
        Hodge-index and height lower bounds.
  \item \emph{Arithmetic B\'ezout local section lengths}
        (\leanlink{GroundTheory.Heights₂.Heights.LocalSectionLengthSum.proof}):
        evaluates the finite-place contribution to arithmetic intersection
        numbers as a sum of finite-length quotients of local section modules on
        an integral model.
\end{itemize}

These height foundations culminate in the unconditional Moriwaki bigness and
arithmetic Hilbert--Samuel theorems in
\leanmod{Development.WeilBound.CoreStatements.AutomorphicKernels.MoriwakiBigness},
which feed directly into the top-level deduction
\leanlink{UpperTheory.LiteratureTheorems.every_character_of_literature_of_vWin}.

\begin{theorem}[Unconditional Moriwaki bigness and arithmetic Hilbert--Samuel]
\label{thm:ch02-moriwaki-big}
\leanlink{PositivityTheory.Pf1.HilbertSamuelMoriwaki.moriwakiBig_outright}
Moriwaki's arithmetic bigness criterion
(\leanlink{PositivityTheory.Pf1.HilbertSamuelMoriwaki.moriwakiBig_outright})
and its arithmetic Hilbert--Samuel application
(\leanlink{PositivityTheory.Pf1.HilbertSamuelMoriwaki.hsUse_outright}) hold
unconditionally for Hermitian line bundles on projective arithmetic varieties:
a Hermitian line bundle with positive arithmetic top self-intersection
$\widehat{c}_1(\overline{L})^{d+1} > 0$ and ample generic fiber is
arithmetically big, admitting exponentially many small integral global
sections in high tensor powers.
\end{theorem}

\begin{proof}
The theorem
\leanlink{PositivityTheory.Pf1.HilbertSamuelMoriwaki.moriwakiBig_outright} is
derived from three unconditionally proved pillars: characteristic-zero
resolution of singularities (\leanlink{PositivityTheory.ResPlug.FromRes.resK}),
the Poincar\'e--Lelong current formula
(\leanlink{PositivityTheory.Tls2.Pl.poincareLelong}), and the resulting
arithmetic intersection system
(\leanlink{GlobalTheory.T.arith_intersection_system_of_poincare_lelong}),
together with the dense-curves Moriwaki reduction
(\leanlink{PositivityTheory.MbAis.DenseCurves.main}); the same inputs
simultaneously establish the arithmetic Hilbert--Samuel application stated
above.
\end{proof}

Finally, the positivity and automorphic-spectrum layers supply the structural
inputs that keep the automorphic representations along the iterated lift chain
well-behaved at both archimedean and finite places:
\begin{enumerate}
  \item \emph{Conjugate-self-dual pullback stability}
        (\leanlink{UpperTheory.CBC.SConjSpaceC_holds}): under a cyclic
        automorphism $\sigma$ of a CM field $E$, pullback preserves the space
        of smooth automorphic forms, intertwines the Lie algebra and center of
        the universal enveloping algebra, and preserves cuspidal
        subrepresentations.
  \item \emph{Generic existence of Satake parameters}
        (\leanlink{UpperTheory.CBC.SatakeParameterExistence.re}): for any
        totally real number field $F$, any rank $n$, and any cuspidal
        subrepresentation $\pi$, there exists a Satake parameter matrix at all
        but finitely many finite places $v$ of $\mathcal{O}_F$.
  \item \emph{Special lattices for nonzero types}
        (\leanlink{UpperTheory.TNZ.step_spLat}): constructs special integral
        lattices witnessing nonvanishing of local automorphic types.
  \item \emph{Theorem of the cube and cyclic-algebra Hilbert~90}
        (\leanlink{GlobalTheory.CompanionTheory.CubePerf.cubeTheorem_of_perfAll},
        \leanlink{GlobalTheory.T.cyclic_algebra_hilbert_ninety}): in the
        abelian-scheme and twisted-factor layers of
        \leanmod{Development.AutomorphicSpectrum}, the theorem of the cube is
        established from spreading of perfect models, and cyclic-algebra
        Hilbert~90 controls local norm-one cocycles on discrete valuation rings.
\end{enumerate}

Together, these results show how an off-line zero of a Dirichlet $L$-function
is transformed---via the critical-line reduction, Selmer windows, Arakelov
lattice slopes, iterated definite theta lifts, integral dual Weyl divisibility,
and Moriwaki arithmetic bigness---into a problem of comparing automorphic periods
and spectral norms.
That comparison is governed by the Jacquet--Rallis relative trace formula and
its accompanying analytic and representation-theoretic estimates, to which we
turn in the next chapter.
